\chapter{Scaling Laws of Systems}
\label{chap:scaling-laws-of-systems}

The apparatus makes its last examination here: how systems scale --- the ceilings, costs, and concentrations
that appear
when many records act together --- the systems ours to build, the ceilings they meet the world's. Six effects
make the case: the cap on parallel speedup, the cost set by the slowest
synchronization, the concentration of mass into a few, agency as strain minimized, the stochastic correction
that is a refinement artifact, and --- the last wall of all --- the entropy that renaming cannot raise. After
them, only the count remains to be read whole.

\section{The Amdahl Effect: a ceiling on parallel speedup}
\label{se:amdahl}

No computation is made arbitrarily fast by parallelism. However many processors are thrown at a problem, the
part that must run in sequence cannot be hurried, and that serial residue caps the gain. If a fraction $p$ of
the work parallelizes and $1 - p$ does not, the speedup is bounded,
\begin{equation}
  S_{\max} = \frac{1}{1 - p},
\end{equation}
no matter how many processors are added: even with infinite parallelism the serial fraction sets the ceiling
--- the processors ours to add, the residue that waits for none of them the world's.
Gene Amdahl drew the bound in 1967 to cool the enthusiasm for ever more processors --- parallelize
nine-tenths of a job and the last tenth, which must run in sequence, still caps the speedup at tenfold, though
a thousand processors stand idle waiting on it.
The honest floor is a finite count obligation --- the speedup is bounded by the serial residue, a fixed
fraction of the work that no parallelism removes --- the splitting ours, the ceiling the world's. The reading
is that parallelism has a ceiling, and it
is the serial part. What cannot be split cannot be hurried, and the gain from any number of processors
is capped at the inverse of the fraction that must run alone.

\section{The Tail Latency Effect: the slowest boundary sets the cost}
\label{se:tail-latency}

The cost of a synchronized system is set by its slowest part. When many processes must reconcile at a shared
boundary, the latency is governed not by the shortest path but by the longest --- the slowest boundary
reconciliation in the system. The tail of the latency distribution scales with the size and the boundary,
\begin{equation}
  \mathcal{L}_{\text{tail}} \propto N\,|\partial\Omega|,
\end{equation}
so that as a system grows, its worst-case synchronization cost grows with it --- the system ours to grow, the
laggard it waits on the world's. The honest floor is a finite
count obligation --- the tail latency is a count of the slowest boundary's work, scaling with the system size
and its boundary --- the boundary ours to set, the tail that finishes last the world's. The reading is that
the slowest boundary sets the cost. A synchronized system waits on its
laggard; the latency is not the average but the tail, the cost of the one reconciliation that finishes last,
and it grows with the boundary the system must keep coherent.

\section{The Pareto Effect: the mass concentrates}
\label{se:pareto}

Refinement concentrates mass into a few. Admissibility seldom spreads weight evenly; extension after
extension, the weight drifts onto a few, so that a small fraction of the entries carry most of the total. The distribution is a
power law,
\begin{equation}
  P(x) \propto x^{-\alpha},
\end{equation}
heavy at the head and long in the tail --- the law Vilfredo Pareto found in the distribution of wealth, the
familiar four-fifths of it held by a fifth of the people, the same lopsidedness repeating within that fifth
and within that again, scale-invariant all the way down --- the fifths ours to draw, the lopsidedness they
find the world's. A sparse basis --- a few entries holding the mass ---
is what refinement keeps returning, while a flat, uniform distribution rarely lasts. The honest floor is a
finite count
obligation --- the concentration is a count, most of the mass on a few entries, the bookkeeping of an observed
lopsidedness rather than a derived necessity --- the spread ours to tally, the concentration it keeps finding
the world's. The reading is that the mass concentrates, and uniformity seldom survives. A
flat distribution is rarely a stable state of a refining system; the weight drifts onto a sparse
few, and the power law is the shape that concentration takes --- the many light, the few heavy.

\section{The Agent Effect: agency is strain minimized}
\label{se:agent}

An agent is not an outside observer but a part of the record that minimizes its own strain. A localized
substructure, facing the stream of events at its boundary, acts to minimize the informational strain those
boundary conditions induce on it --- and what we call agency is the residue of that minimization, not a force
added from outside. The agent is the part of the record that pushes back on its own boundary to keep its
strain low --- the ledger ours to reconcile, the demand the stream sets the world's. The honest floor is a
finite count obligation --- agency is the strain-minimizing residue, a count
of how a local substructure reconciles its boundary; the agent counts its refinements, and nothing here
selects them --- the reconciliation ours, the strain that remains the world's. The reading is that agency is
strain minimized, not a
special cause. An agent is a region of the record that keeps its own informational strain low against the
incoming stream, and its apparent purpose is the shape that minimization takes --- not a force outside the
physics but a residue within it.

\section{The Ito Lemma: a stochastic correction that is a shadow}
\label{se:ito}

When a record is updated under noise, a correction appears that is an artifact of refinement, not an extra
force. Step a noisy variable forward, $X_{t+\Delta t} = X_t + \Delta X_t$, and ask how a function of it
changes; the answer carries a second-order term the ordinary chain rule would miss,
\begin{equation}
  df = f'(X_t)\,dX_t + \tfrac{1}{2} f''(X_t)\,(dX_t)^2,
\end{equation}
where the $(dX_t)^2$ term --- the quadratic variation --- does not vanish, because under noise the squared
step is of the same order as the step. For the Brownian walk already met, the rule is exact, $(dX_t)^2
\to dt$: the variance of a random step grows with the time elapsed, not with its square, so the squared step
keeps pace with the step itself and the second-order term refuses to disappear --- the function ours to write,
the correction the jitter exacts the world's. This is Ito's lemma, and the correction is the smooth shadow of a
discrete reconciliation: it is forced by comparing the refined updates, not by any force in the dynamics. The
honest floor is a smooth-shadow analogy --- the discrete noisy updates are the model, the Ito correction their
smooth shadow, and the bridge from the squared discrete step to the continuum quadratic-variation term is an
analogy, named as one --- the calculus ours to bring, the quadratic variation the world's. The reading is that
the stochastic correction is a refinement artifact, not a force.
The extra term in Ito's lemma is what the discrete reconciliation leaves when it is read in the continuum ---
a shadow of how the noisy steps compare, and the last smooth shadow of them all.

\section{The Limitation of Indexing: the entropy renaming cannot raise}
\label{se:limitation-indexing}

The last wall has been standing since the first mark. Entropy --- the logarithm of the number of
distinguishable configurations,
\begin{equation}
  S = \ln N, \qquad \Delta S \ge 0
\end{equation}
--- rises only under a genuine refinement, a real new distinction committed to the record. It does not rise
under renaming. Relabel the entries, permute the indices, rewrite the record in another alphabet, and $N$ is
unchanged and $\Delta S$ is zero --- the relabeling ours, the count it cannot move the world's: a renaming
adds no distinguishable configuration, and so adds no entropy.
Only a real new distinction --- a measurement that separates two configurations the record could not tell apart
before --- raises the count. The honest floor is a no-go, the last of all: the count cannot be
inflated by renaming, only by measuring --- the labels ours, the count the world's. The reading is the first
difference, restated as a wall. The
first mark made a difference because it distinguished two configurations; a relabeling makes none, and so the
entropy a record carries is exactly the count of real distinctions in it, never the count of names. Renaming
measures nothing. The arrow that began at the first difference points the same way at the last: the count
rises only when something is genuinely told apart.

\bigskip
\begin{center}
\rule{0.4\textwidth}{0.4pt}
\end{center}
\bigskip

\noindent\textbf{\large An Entire Physics, Read Off the Count}

\medskip

This book began with a single difference. In its first chapter a mark was made, and the mark mattered because
it told two configurations apart --- the first distinction, the first count. Everything after was the same act,
repeated and refined: a count incremented, an order kept, a distinction committed and never retracted. From
that one difference, an entire physics has been read off the count.

And it was read by one instrument. Not a different apparatus for each domain --- a telescope here, a cloud
chamber there, a clock, a lattice, a computer --- but one finite reader, the same from the first mark to the
last: a thing of finite resolution, finite memory, and a single primitive act, the recording of a difference.
Everything in this book was that one instrument at work, carried from mechanics to the horizon and back, and
the reader who held it in mind from the first count will have recognized it behind the last.

The path ran through six parts. Mechanics gave the count its invariance under relabeling and its least
principle. Fields turned the count into a holonomy that survives where the field is absent and a frontier that
dilutes as it grows. Heat made the count a temperature and proved that information is physical --- that to
sort, to erase, to confine is to pay in the currency of the second law. The quantum carried the count as an
amplitude, charged a loop to plus or minus one, and --- at the book's deepest grounding --- registered a
particle at a detector, the positron proved by a coincidence the model could state as a theorem. Relativity
read the interval, the potential, and the horizon as the geometry of the count, and forbade the one closed
loop in time the count's own rising will not allow. And the foundations turned the apparatus on itself, finding
the preconditions it cannot escape, the continuum it casts as a shadow, the walls of computation it cannot
pass, and the entropy that renaming cannot raise.

Through all of it, every effect was a count, and every count was fenced honestly at one of four ceilings. Some
were finite count obligations --- a balance, a conservation, a closure the record owes. Some were smooth-shadow
analogies --- a continuum law named as the shadow the discrete casts as it refines, the bridge to the equation
called an analogy and not a derivation. Some were no-go walls --- a refinement the law makes impossible, from
the first friction threshold to the last forbidden renaming. And some were names only --- a phenomenon written
in full as physics, with the finite model claiming nothing but the label beneath it. The discipline was to say,
of every effect, exactly which floor it stood on, and never to claim more.

A handful stood higher than the rest. In a few places the finite model did not shadow physics but proved it:
the echo that charges a closed loop, the Bragg peak that proves a matter wave, the Dirac operator and the
chirality and the Feynman loop that carry a proved $\pm 1$, the Sagnac residue that will not close, and, above
all, the positron --- detected, by a coincidence the model proves as an if-and-only-if, with every converse
falsifiable, and fenced with equal care against claiming one inch more than it had earned. Those are the places
where the book's central wager was paid in full: that a finite record, kept honestly, can not only describe the
physics but, here and there, prove it.

One scene remains, and the whole of the physics sits in it. A driver is clocked for speeding, and three
instruments say so. The speedometer is the car's own count --- an origin fixed, the distance ticked off, the
rate read from inside the motion: the first regime of this book, a record keeping the books on itself. The
GPS is the second, and it reads true only because relativity is true: the satellite clocks run fast by the
thirty-eight microseconds a day the potential wins --- the same clock physics Pound and Rebka weighed in a
Harvard stairwell --- and uncorrected, the position, and the speed read off it, drift by kilometers within
hours. The radar gun is the third: a photon sent out, reflected off the car, returned with its frequency
shifted, the shift converted to a speed --- the quantum doing the measuring, the frame physics doing the
converting. Three instruments, three regimes; one car, and a court that owes one verdict.

The court's first work is to close the three readings into one. Each instrument went out by its own physics
and came back with a number, and the comparison of the returns is a closed loop --- the same figure the
fields and the quantum charged: what survives when the circuit is sewn shut is a residue no re-reading
cancels. When the three numbers agree, the agreement is a conservation, residue composing with residue along
the loop; when they disagree, the disagreement is itself a datum, a residue pointing at whichever leg of the
physics bent. Either way the loop returns something the court can hold: one speed, carried by the concord of
three independent ways of counting it.

Then the harder question: whose speed is it? A speed is a reading against a frame --- the road's, in which
the car moves, or the car's, in which the road does --- and the world does not fix which frame the question
is asked in. That is not a defect in the instruments; it is the frame physics doing here what it did on its
own pages. The sign a turned loop hands back is the frame's, not the world's; and which particle a detector
calls matter is the apparatus's to say --- the coincidence that registered the positron is written in the
two photons alone, and no line of it names which partner was the matter. Every reading on
the bench is sound, and not one of them can supply the frame it is to be read in.

And so the question itself turns. The case opened on a question put to the world --- \emph{is he speeding?}
--- a fact, to be measured. It closes on a question put to a juror --- \emph{do you find him guilty?} --- a
verdict, to be returned. The same case, the same car, the same number: the question has made one full
revolution, from the world's frame to the asking's, and come back wearing the same face with the opposite
sign. The electron whose sign Part IV followed is a spin-half thing --- a spinor --- and a spinor turned
through one full revolution, $2\pi$, does not come home: it returns with a minus sign, the same face
inverted, and only a second full turn, $4\pi$, restores it. That is the case's own geometry. The opening
question and the closing one are the $2\pi$ pair --- fact and verdict, one face with two signs.

But before the verdict, the floor. The court does not prove the speed; it stands on the instruments, and
the instruments stand on the physics. Trace the radar down: the photon's emission and its shifted return
rest on the quantum; the conversion of shift to speed rests on the frames --- rests on, in each case a
dependency and never a derivation, the same fence this book kept everywhere else. And relativity itself is
read by a clock --- the thirty-eight microseconds are a clock's gain --- and the clock keeps an assumption
older than either theory: Newton's, that time is a continuous parameter, a smooth $t$ that flows. Einstein
localized the clock, relativized it, and kept the $t$ continuous. And the continuity of time rests on
nothing at all. It has never been observed: what is observed is ticks, a finite count, never the continuum
between them. It has never been derived. It is the assumption adopted for want of anything else to do ---
and this book has been, from the first mark, the something else: a physics run on the count alone, the
continuum fenced, wherever it appeared, as the smooth shadow the count casts. Not disproved --- the fences
hold, and no more is claimed at the end than anywhere earlier --- but named, and named as an assumption.

So the court convicts, if it convicts, standing on a plank it cannot prove --- and the juror it turns to has
been seated since the first mark. Whoever held the instrument through these chapters --- watched the count
rise, the fences hold, the continuum recede into shadow --- was never watching from the gallery. The case
closes on its second question, the one no instrument answers: which frame, which sign, guilty or not. The
reading of the readings was never the apparatus's to give. It is the juror's --- and the juror is you. One
revolution brought the question back with its sign inverted; the second turn is yours to make. The
verdict is yours.

\bigskip

\hfill Q.E.D.
